\subsection{Generator Evaluation Protocol}
\label{sec:generator-evaluation-protocol}

\subsubsection{Autoregressive Approximation Studies}
\label{sec:ar-coverage-protocol}
The executed KernelSynth prediction study evaluates 33 entries of the official
kernel bank and 20 sampled compositions. For each mechanism it draws 16 training
and 16 independent test trajectories of length 1024. A pooled linear AR model
with an intercept is fitted by ridge regression, with penalty $10^{-5}$ on
the lag coefficients. The orders tested are
$\{1,2,4,8,16,32,64,128,256\}$. Predictions use the observed preceding values
at every test position. The reported score is the pooled mean squared error
divided by the pooled test-target variance plus $10^{-12}$.

The companion TimePFN experiment evaluates 33 latent series from 11 LMC
mechanisms, using three latents per mechanism and orders up to 128. This is an
\emph{oracle-latent} experiment: the latent series are supplied directly to
the AR models. Order selection uses the smallest score on the same test set
being summarized, so the results describe exploratory best-of-grid
approximation capacity. The saved results and links to the executed records
appear in Section~\ref{sec:ar-coverage-results}.

\subsubsection{Four-Family Stability Comparison}
The executed notebook
\path{experiments/abstract_syntax_tree/6_stability_methods_tradeoff.ipynb}
uses root seed 20261001, 5,000 independent mechanisms per family, and one
trajectory per mechanism. Each run has one channel, a lag range of 1--12,
and diagnostic horizons $h\in\{96,192,384\}$. The original trees contain
1--6 binary operations, addition probability 0.8, parameter/lag leaf
probabilities 0.55/0.45, and coefficient magnitudes in $[0.05,0.8]$ with
random signs. Explicit time leaves are disabled in this comparison.

The evaluator applies the explosion and tail diagnostics defined in
Section~\ref{sec:generator-diagnostics}. It retains every attempt in the
denominators and reuses the prefixes of one trajectory for the different
horizons. The family-specific construction, initialization, and innovation
laws are described in Section~\ref{sec:generator-stability}; results identify
the historical all-lag Shao--Wu saturation policy used in this completed
comparison.

\subsubsection{Timing and Repeated-Run Protocol}
Timing includes sampling/certification, initialization, innovation generation,
AST execution, and diagnostics, including the tail measurements and their
three-tolerance sensitivity analysis. Completed full-horizon trajectories per
second uses the total wall-clock cost, including failed attempts. This prevents
early-stopping generators from appearing efficient solely because they
evaluate fewer time steps.

For timing uncertainty, Notebook~6 executes the complete 5,000-mechanism
workflow five times per family on the same machine, with BLAS restricted to
one thread. Each run rebuilds the mechanisms with the same seeds, samples
histories and innovations, and evaluates the same diagnostics. The method
order rotates in the additional runs. Identical seeds make the outcome counts
reproducible, so the spread in wall-clock time reflects execution conditions
rather than a different draw of mechanisms. We report the mean and sample
standard deviation across runs, plus their observed minimum--maximum range.
These are descriptive measurements, not confidence intervals for performance
on other hardware. The first run is included without a separate cache warm-up.
